\scanpage{12}
45bis

\margin{\textcircled{x}}
(v)- Supposons $\underline{a}$ rationnel sur $K$, et $\underline{q}$ définie sur $K$. Alors $\xi_0(\underline{q}_{\underline{a}_m})$ est définie sur $K$ pour tout $m$. Soit $x$ un point de $A$, rationnel sur une extension algébrique de degré fini $L$ de $K$. Pour $\underline{q}$, $\underline{a}$, $x$ et $m$ fixés, le nombre $(\xi_0(\underline{q}_{\underline{a}_m}))((x)-(0),v)$ ne dépend que de la valeur absolue $w$ induite par $v$ sur $L$. \margin{\figure{\begin{tikzpicture}\draw[thick] (0,0)--(.03,-.35);\end{tikzpicture}}}Par passage à \deleted{\unsure{l'infini}} la limite pour $m$ infini, on en déduit qu'il en est de même de $(\mu(\underline{q},\underline{a}))((x)-(0),v)$, autrement dit que $\mu(\underline{q},\underline{a})$ est définie sur $K$.
\margin{$\longrightarrow$ \textcircled{xx} p.45}
\mspage{11}
